\section{Versions A and B compared}\label{sec:sp:compare}\label{sec:sp:B:compare}

Sections~\ref{sec:sp:A} and~\ref{sec:sp:B} solve the same planner's problem under two closures of one
object, the material intensity of new capital goods. Version A ties it to the economy-wide intensity,
$\Phi^I=\Omega\phi^I$; Version B takes it exogenous. Each section was written to be read on its own,
and neither refers to the other. This section does nothing else: it states exactly how the single
assumption propagates, how large the difference is, what survives it untouched, and what turns on the
choice. Nothing here is a robustness check --- the two closures answer different questions about what
capital is made of, and the differences are structural rather than quantitative.


\subsection{The whole difference is one term in the Lagrangian}\label{sec:sp:compare:map}

Start from the constraint sets. Substituting $\Phi^I=\Omega\phi^I$ into Problem~\ref{prob:sp:B} makes
its ledger~\eqref{eq:sp:B:ledger} and its law of motion for $M^K$ identical to
\eqref{eq:sp:A:ledger} and its counterpart in~\eqref{eq:sp:A:states-motion}. What remains is that
Version A adjoins $\Omega$ as a control together with the intensity
definition~\eqref{eq:sp:A:intensity}, and drops the mass inequality~\eqref{eq:sp:B:mass}, which is
vacuous once $\Omega$ adjusts. So the two problems differ in exactly one place, and it is visible in
the Lagrangians. Writing $\lambda\zeta$ for the multiplier on~\eqref{eq:sp:A:intensity} and using
$R=N+R^R$,
\begin{align}
\underbrace{\lambda\zeta\Big[N+R^R-\Omega\left(\mathcal D+\phi^IG\right)\Big]}_{\text{Version A}}
\;=\;
\underbrace{\lambda\zeta\Big[N+R^R-\Phi^IG\Big]}_{\text{Version B's term, with }\eta\mapsto\zeta}
\;-\;\lambda\Delta\mathcal D,
\qquad
\Delta\equiv\zeta\Omega .
\label{eq:sp:compare:lagdiff}
\end{align}
Everything that distinguishes the two systems of optimality conditions is generated by the single
extra term $-\lambda\Delta\mathcal D$, together with the change in status of the multiplier on
$R-\Phi^IG$ from an inequality multiplier to an equality multiplier. Recall from~\eqref{eq:sp:A:zeta}
the objects it is built from,
\begin{align}
\chi &\equiv \frac{\Phi^IG}{R}\in\left[0,1\right),
&
\zeta &= \chi\left(p^W-p^M\right),
&
\Delta &= \zeta\Omega ,
\label{eq:sp:B:compare:Delta}
\end{align}
the share of material input embodied in new capital, and the wedge it generates.

Because $\mathcal D=\omega^YY+\omega^NC^N+\omega^DC^D+\omega^WC^W+\omega^CC$ is linear in the
final-good flows, differentiating $-\lambda\Delta\mathcal D$ with respect to any control other than
$\Omega$ multiplies each flow's own coefficient by $\Delta$ and nothing else. That gives the
correspondence in closed form.

\begin{proposition}[Exact correspondence of the two systems]\label{prop:sp:compare:map}
Let $\Delta$ and $\zeta$ be as in~\eqref{eq:sp:B:compare:Delta}. Every first-order and costate
condition of Problem~\ref{prob:sp:A} is the corresponding condition of Problem~\ref{prob:sp:B} under
the substitutions
\begin{enumerate}[label=(\roman*),leftmargin=2.2em,itemsep=0.15em]
\item $\Phi^I\mapsto\Omega\phi^I$ throughout;
\item $\eta\mapsto\zeta$, the multiplier on the mass inequality replaced by the multiplier on the
intensity definition;
\item every appearance of $F$ and its derivatives $F_K,F_R,F_P$ --- including those inside the
composites $B$ and $d$ --- multiplied by $\left(1-\Delta\omega^Y\right)$;
\item every appearance of an outlay of the final good and its derivatives multiplied by
$\left(1+\Delta\omega^j\right)$: $C^N_N$ and $C^N_S$ by $\left(1+\Delta\omega^N\right)$, $C^D_D$ and
$C^D_X$ by $\left(1+\Delta\omega^D\right)$, $c^W$ and $dc^T/d\varpi$ by
$\left(1+\Delta\omega^W\right)$, and the consumption condition $u'(C)=\lambda$ by
$u'(C)=\lambda\left(1+\Delta\omega^C\right)$.
\end{enumerate}
Gross capital formation is not repriced: $G$ enters the base with weight $\phi^I$ rather than an
$\omega^I$, and its repricing appears instead as the factor $\left(1-\chi\right)$
in~\eqref{eq:sp:A:foc:I}. Setting $\Delta=\zeta=0$ and reinstating $\eta$ returns
Section~\ref{sec:sp:B:foc} exactly.
\end{proposition}

\begin{proof}
By~\eqref{eq:sp:compare:lagdiff} the two Lagrangians differ by $-\lambda\Delta\mathcal D$ once
$\Phi^I=\Omega\phi^I$ and $\eta$ is renamed $\zeta$. Adding $-\lambda\Delta\mathcal D$ to the goods
term of~\eqref{eq:sp:B:lagrangian} collects to
\begin{align}
\lambda\Big[\left(1-\Delta\omega^Y\right)F-\left(1+\Delta\omega^C\right)C-I-\delta K
-\left(1+\Delta\omega^N\right)C^N-\left(1+\Delta\omega^D\right)C^D
-\left(1+\Delta\omega^W\right)c^W\!\left(\varpi\right)W\Big],
\notag
\end{align}
since $Y=F$ and $C^W=c^W(\varpi)W$, and since $I+\delta K$ appears in no term of $\mathcal D$.
Differentiating this expression with respect to any control or state reproduces~(iii) and~(iv), the
chain rule carrying the factor $\left(1-\Delta\omega^Y\right)$ through every route by which a control
moves $Y$ --- which is why it appears inside $B$ and inside the damage $d$ as well as in the
conditions for $K^R$ and $K$. For~(ii) and the investment claim, the terms in $I$
are $-\lambda+\lambda q-\lambda\Omega\phi^I\left(p^M-p^W+\zeta\right)=0$ under A and
$-\lambda+\lambda q-\lambda\Phi^I\left(p^M-p^W+\eta\right)=0$ under B, so the two coincide
under~(i)--(ii); substituting $\zeta=\chi\left(p^W-p^M\right)$ into the first gives
$q=1-\Phi^I\left(1-\chi\right)\left(p^W-p^M\right)$, which is~\eqref{eq:sp:A:foc:I}.
\end{proof}

Three things about this proposition are worth being careful about, because a loose reading of it
would suggest more than it says.

First, it is a correspondence between \emph{systems of conditions}, not a change of variables that
solves Version A. $\Delta$, $\zeta$, $\chi$ and $\Omega$ are endogenous, and $\Delta$ depends on
$p^W-p^M$, which the conditions themselves determine. The map is therefore evaluated at the Version-A
allocation, and it identifies the form of the difference rather than eliminating it.

Second, the two multipliers exchanged in~(ii) are not the same kind of object, and this is where the
economic difference sits. $\eta\ge0$ obeys complementary slackness and is \emph{zero except at a
corner} of the accounting; $\zeta$ is an equality multiplier, free in sign and generically nonzero.
The substitution therefore replaces a generically inactive shadow price with a generically active
one. That single change of status, not the algebra of the repricing, is why Version A carries wedges
everywhere the final good is spent and Version B carries none.

Third, the correspondence is silent about the solutions. Two problems whose conditions map into each
other under an endogenous substitution need not have allocations that are close, and
Section~\ref{sec:sp:compare:nesting} takes this up.


\subsection{How large the difference is}\label{sec:sp:compare:size}

The wedge admits a reading in observable units. Since $\Delta=\Omega\chi\left(p^W-p^M\right)$, the
wedge attached to activity $j$ is
\begin{align}
\Delta\omega^j &= \underbrace{\Omega^j}_{\text{tonnes drawn per unit spent on }j}
\times\underbrace{\chi}_{\text{share of throughput stored in durables}}
\times\underbrace{\left(p^W-p^M\right)}_{\text{value of postponing a tonne}},
\qquad
\Omega^j\equiv\Omega\omega^j,
\label{eq:sp:compare:size}
\end{align}
a product of three quantities each of which has a direct empirical counterpart. $\Omega^j$ is the
absolute material intensity of activity $j$, the object material-flow accounts report; $\chi$ is the
share of material input embodied in new capital; and $p^W-p^M$ is the difference between shedding a
tonne now and shedding it over the life of an asset. Along a path on which $p^W$ is constant,
\eqref{eq:sp:A:costate:M} gives $p^M=\delta p^W/\left(r+\delta\right)$ and hence
\begin{align}
\Delta\omega^j &= \Omega^j\chi\,p^W\,\frac{r}{r+\delta},
\label{eq:sp:compare:size-stationary}
\end{align}
so the difference between the closures is second order in exactly the circumstances one would expect:
it vanishes as $r\to0$, when postponement is worthless; it vanishes as $\delta\to\infty$, when capital
stores nothing; and it is proportional to $\chi$, which is small in an economy that embodies little of
its throughput in durables. It is worth reporting $\chi$ alongside $a\varpi$ for this reason ---
$\chi$ is the statistic that says how much the choice of closure can matter, and
\eqref{eq:sp:compare:size-stationary} converts it into a number.

Two qualifications. The formula is a level statement, not a bound on the difference in allocations: a
small wedge on every margin simultaneously can still move a transition date, since the transition is
driven by a slow-moving rent and small persistent wedges accumulate. And $\chi$ need not become small
in the long run. As the economy approaches the circular boundary, $N\to0$ with $R$ bounded, so
$\chi=\Phi^IG/R$ need not vanish and neither need $\Delta$: the two closures do not converge as the
economy becomes circular, which is the region the paper is about.


\subsection{What the closure does not reach}\label{sec:sp:compare:common}

Four results hold under both closures with the same proof, and they are the results the paper turns
on.

\smalltitle{The materials arbitrage.} $B^{\mathrm B}=p^{N,\mathrm B}+d^Wp^P$
in~\eqref{eq:sp:B:arbitrage} and $B=p^N+d^Wp^P$ in~\eqref{eq:sp:A:arbitrage} are the same statement
with the same cancellation: the shadow price of waste drops out of the comparison of sources under
either closure. Recycling defers waste rather than avoiding it, and no setting of the parameters
recovers a specification in which it avoids waste --- mass conservation is not a modelling option in
this framework. The companion statements \eqref{eq:sp:B:FR} and~\eqref{eq:sp:A:FR} likewise coincide
in form, with $p^W$ charged exactly once, at entry.

\smalltitle{Scarcity raises circularity.} Propositions~\ref{prop:sp:B:circularity}
and~\ref{prop:sp:A:circularity} are proved identically, from the curvature of $a$ and the convexity of
$c^T$ alone. One caveat, which the parallel statement obscures: the two propositions hold different
things fixed. Under B the ceteris paribus is $F_K$ and $p^P$; under A it is
$\left(1-\Delta\omega^Y\right)F_K$, $\left(1+\Delta\omega^W\right)$ and $p^P$. The mechanism is
identical; what the wedge does is move the terms held fixed, which is a second-order effect on the
same margin.

\smalltitle{No interior steady state.} Propositions~\ref{prop:sp:B:nosteadystate}
and~\ref{prop:sp:A:nosteadystate} use only the state equations and the ledger, both common to the two
closures. The obstruction is $\dot X=D$, not the treatment of $\Phi^I$, and the repair $C^D_X=0$ is
the same under both.

\smalltitle{Capital is a temporary materials sink.} $q$ lies below its goods cost whenever $p^W>0$
under either closure, and the sign flip of $p^W$, $p^M$ and the capital wedge is simultaneous under
both. Version A adds $\Delta$ to the list of objects that flip together, by~\eqref{eq:sp:A:zeta}.

To this can be added a fact that matters for the quantitative implementation: the two closures carry
the \emph{same} states and the same shadow prices, $q,p^S,p^X,p^P,p^W,p^M$. Neither adds a state
variable or a price to the other. Only the residual equations differ, and only through $\Delta$,
$\zeta$ and the exchange of the complementarity block for the intensity definition.


\subsection{Where they differ}\label{sec:sp:compare:table}

Table~\ref{tab:sp:versions} collects the differences.

\begin{table}[H]
\centering
\caption{Version B and Version A compared.}
\label{tab:sp:versions}
\setlength{\tabcolsep}{5pt}
\renewcommand{\arraystretch}{1.05}
\begin{tabular}{@{}p{0.26\textwidth}p{0.33\textwidth}p{0.33\textwidth}@{}}
\toprule
 & Version B ($\Phi^I$ exogenous) & Version A ($\Phi^I=\Omega\phi^I$) \\
\midrule
Materials restriction & $\Phi^IG\le R$ binds as an inequality, multiplier $\eta\ge0$
& none; $\Omega\ge0$ holds by construction \\
Status of that multiplier & zero except at a corner & free in sign, generically nonzero \\
Solving for $W$ & one division, \eqref{eq:sp:setup:ledger-solved}
& positive root of a quadratic, \eqref{eq:sp:setup:ledger-quadratic} \\
$\Omega$ & computed after the fact & solved jointly, a control \\
Relative intensities $\omega^j$ & do not affect the allocation & affect every condition \\
Wedge on spending & none & $1+\Delta\omega^j$ \\
Wedge on output & none & $1-\Delta\omega^Y$ \\
$\lambda$ & equals $u'(C)$ & equals $u'(C)/\left(1+\Delta\omega^C\right)$ \\
Euler equation & standard, \eqref{eq:sp:B:euler} & extra term in $\frac{d}{dt}\ln\left(1+\Delta\omega^C\right)$ \\
Capital price $q$ & $1-\Phi^I\left(p^W-p^M\right)$ & $1-\Phi^I\left(1-\chi\right)\left(p^W-p^M\right)$ \\
$p^W$ & explicit, \eqref{eq:sp:B:foc:W-solved} & implicit, \eqref{eq:sp:A:foc:W-solved} \\
Signs that flip together & $p^W,p^M,q$ & $p^W,p^M,q,\Delta$ \\
Circular limit & $\Delta$ absent throughout & $\chi$, and hence $\Delta$, need not vanish \\
Planner's information & needs no $\omega^j$ & needs every $\omega^j$ \\
\bottomrule
\end{tabular}
\end{table}

The row on $p^W$ deserves a word, since it is the one difference that is felt immediately in any
numerical work. Under B, \eqref{eq:sp:B:foc:W-solved} is a formula: given extraction costs, rents,
damages and the recycling technology, $p^W$ is evaluated. Under A the same equation contains $\Delta$,
which by~\eqref{eq:sp:B:compare:Delta} contains $p^W-p^M$, so $p^W$ solves a fixed point at every
date. Nothing establishes that the fixed point is unique or that the natural iteration converges; both
are open, and both belong with the sufficiency question of Remarks~\ref{rem:sp:A:soc}
and~\ref{rem:sp:B:soc}.


\subsection{The genuine wedge and the spurious one}\label{sec:sp:compare:wedges}

The comparison worth dwelling on is between Version A's consumption wedge and the artefact of
Section~\ref{sec:sp:B:nowedge}. The two have the same shape.
\eqref{eq:sp:B:spurious-C} gives $u'(C)=\lambda\left(1+\Omega^Cp^W\right)$ from a posited reduced form
with fixed intensities; \eqref{eq:sp:A:foc:C} gives $u'(C)=\lambda\left(1+\Delta\omega^C\right)$ from
Version A. Both look like a consumption tax levied on environmental grounds, and only one of them is
real.

They are told apart by the multiplier, and by~\eqref{eq:sp:compare:size} the comparison is exact: the
spurious wedge is $\Omega^Cp^W$, the genuine one is $\Omega^C\chi\left(p^W-p^M\right)$. The two differ
by the factor $\chi\left(p^W-p^M\right)/p^W$, which is at most $\chi<1$ and vanishes with $r$. That
factor is the whole content of the distinction. Under the reduced form, consuming a unit less reduces
economy-wide waste without reducing the mass drawn in, which conservation of mass does not permit ---
the wedge prices a margin that does not exist, and it prices it at the full value of a tonne of waste.
Under Version A total mass is still conserved, and consuming a unit less shifts a given mass out of
the current waste flow into the durable stock; the wedge prices the \emph{timing} of disposal, not its
quantity, which is why it collapses to zero when timing is worthless ($p^W=p^M$) or when nothing is
stored ($\chi=0$). Version B has no wedge at all because it has no such margin either: there $\Phi^I$
is fixed, so the split between stock and flow is not something the composition of spending can move.

The practical statement is that a reduced-form model overstates the case for a consumption tax by the
factor $p^W/\left[\chi\left(p^W-p^M\right)\right]$, which is large whenever $\chi$ is small --- that
is, whenever the economy embodies little of its throughput in durables, which is the empirically
relevant case.


\subsection{Nesting, and what would discriminate}\label{sec:sp:compare:nesting}

Neither closure is a special case of the other, and it is worth being exact about why, since
Proposition~\ref{prop:sp:compare:map} might suggest otherwise.

Suppose one takes the Version-A solution and feeds its intensity path into Version B, setting the
exogenous $\Phi^I(t):=\Omega(t)\phi^I$. The two problems then have the same constraint set along that
path. They still do not have the same solution, because the Version-B planner holds $\Phi^I$ fixed
when perturbing the composition of spending while the Version-A planner does not; the difference is
the feedback $\partial\Phi^I/\partial\left(\text{spending}\right)$, which is precisely the term
$-\lambda\Delta\mathcal D$ of~\eqref{eq:sp:compare:lagdiff}. The A allocation is therefore in general
not optimal for the B problem constructed to agree with it, and conversely. The two families of
problems intersect in one point only: $\Phi^I\equiv0$, that is $\phi^I=0$ under A and $\Phi^I=0$ under
B, where the feedback term is identically zero and the mass inequality never binds.

\begin{remark}[When the two versions coincide]\label{rem:sp:coincide}
$\Delta=0$ iff $\chi\left(p^W-p^M\right)=0$, and the three ways this can happen are not on the same
footing. $\phi^I=0$ is a restriction on the primitives and gives coincidence of the two \emph{problems},
as just described --- this is the minimal model of Section~\ref{sec:sp:schemes}. $G=0$ and $p^W=p^M$ are
properties of an allocation, not assumptions one can impose: they deliver coincidence of the
\emph{conditions} at the dates where they hold and nothing more. Since $p^W=p^M$ requires $r=0$ or
$p^W=0$ by~\eqref{eq:sp:A:costate:M}, it is a knife-edge rather than a case.
\end{remark}

That the intersection is $\Phi^I\equiv0$ has a convenient consequence for the quantitative work. The
stripped model with $\Phi^I=0$ that Section~\ref{sec:sp:schemes} identifies as the natural first
numerical target is closure-free: it is simultaneously Version A with $\phi^I=0$ and Version B with
$\Phi^I=0$, so nothing computed in it depends on the fork, and the fork can be deferred until the
embodied-material block is switched on.

\smalltitle{What would discriminate.} The two closures have a sharp observational difference, and it
is not about mass flows. Under Version B the allocation is invariant to the relative intensities
(Proposition~\ref{prop:sp:B:irrelevance}), so two economies with identical technologies, endowments
and preferences over aggregates, but different \emph{compositions} of final-good spending, have
identical allocations and differ only in reported material intensity. Under Version A they do not:
by~\eqref{eq:sp:compare:size} a composition that raises $\Omega^C$ raises the wedge on consumption,
and changes the recycling rate, the capital price and the transition date. A shift in the composition
of final demand that leaves aggregate throughput unchanged is therefore allocationally inert under B
and consequential under A. Whether that difference is detectable is a separate question --- the
wedges are of order $\chi p^W r/(r+\delta)$ --- but it is the difference the data would have to speak
to, and it is a difference in the response of the allocation, not in the accounts.


\subsection{What the choice of closure decides}\label{sec:sp:compare:decides}

Three things, and none of them is a robustness check.

\smalltitle{What the calibration exercise is.} Under Version B the relative intensities are a
measurement bridge: they translate the model's mass flows into published material-flow accounts and do
nothing else, and by Proposition~\ref{prop:sp:B:irrelevance} no observation of the allocation
identifies them. Under Version A they are technology, identified in principle from the allocation
itself, and they must be estimated rather than merely read off the accounts. The calibration strategy
therefore has to be settled jointly with the choice of closure, not after it;
see~\ref{op:sp:identify}.

\smalltitle{How much of the process-level accounting is load-bearing.} Under Version B
Section~\ref{sec:sp:setup:materialaccounting} can be read as bookkeeping, and the process-level detail
earns its keep entirely through what it \emph{removes} --- the spurious wedges of
Section~\ref{sec:sp:compare:wedges}. Under Version A the same detail is part of the technology and
determines the allocation directly. The methodological claim of the paper is therefore stronger under
A and cleaner under B.

\smalltitle{What the government has to know.} Under B no optimal instrument depends on any $\omega^j$
or on $\Omega$: the least measurable part of the accounting is also the part the policymaker does not
need. Under A every instrument does. This is developed in Part~\ref{part:mkt}, and it is the practical
form the fork takes.

Neither closure is adopted as the baseline here. Both are carried in full, and
Section~\ref{sec:sp:schemes} classifies the available simplifications under both.
